\begin{abstract}
Security research on large language model (LLM) agents is bottlenecked by its
instruments: the strongest existing benchmarks---Agent Security Bench (ASB),
AgentDojo, InjecAgent, and the ToolEmu sandbox---are single-agent, and their only
cross-agent vector is a shared memory store. Yet deployed systems are increasingly
orchestrated collectives in which a planner directs workers, agents share memory,
tools are third-party, and coordination is forced. The phenomena that dominate risk
in that regime---collusion among agents, a compromised orchestrator, malicious
tools, Sybil identities, and group-level emergent failure---cannot be produced, let
alone measured, by a single-agent harness. We present \defname, an environment and
evaluation methodology purpose-built for the security of multi-agent LLM systems.
\defname{} contributes: (i) a task suite with heterogeneous roles (orchestrator,
workers, tool-executors, shared-memory store), partial observability, and
\emph{forced-coordination} tasks that no single agent can complete alone, over star,
chain, tree, and mesh topologies at $3$--$50$ agents; (ii) a natively parameterized
adversary model spanning adversary fraction, independent-versus-colluding behaviour,
static-versus-adaptive strategy, orchestrator compromise, malicious tools, Sybils,
and poisoned shared memory; (iii) a system-level metric suite, including a
generalization $\NRPs=\PNAs\cdot(1-\ASRs)$ of ASB's utility--security metric that is
scored against deterministic security predicates over the global state trace, never
an LLM judge on the critical path, so a successful injection cannot hijack the
evaluator; and (iv) \corb, a co-evolving red/blue protocol that alternates attack and
defense proposals and logs their system-level scores. We develop the theory the
instrument requires rather than assuming it: an axiomatic characterization that
singles out the product aggregator and reduces it to ASB's metric at one agent;
an evaluator-integrity theorem separating deterministic checking from LLM-judge
hijack; topology-dependent compromise-propagation bounds and a
collusion-versus-independent separation in the environment's cascade model; a
forced-coordination lower bound; and finite-sample concentration for the estimators.
Built on ToolEmu-style LM emulation (human agreement $\kappaem\approx 0.48$, which we
surface as a first-class caveat) and AgentDojo-style deterministic checks, \defname{}
is an instrument, not a verdict: we never equate a lower system-level attack success
rate with ``secure.'' Because the backbone interface is model-agnostic, we evaluate the
instrument end-to-end on a transparent, clearly-labelled \emph{synthetic} agent emulator
(all headlined output in this version is synthetic and provenance-stamped; a real-backbone
pilot of six open instruct models over twelve configurations is reported separately and
explicitly labelled PILOT). On that emulator the headline claim is measured: defenses effective in
the single-agent regime lose advantage under collusion and compromised-orchestrator
conditions (a per-agent guardrail's relative ASR reduction over no-defense shrinks
$41.6\%\to31.8\%\to27.6\%$ across those contexts; Table~\ref{tab:degradation}); replication
on real backbones and the deferred topologies/scales is the release-phase programme.
\end{abstract}

\begin{highlights}
\item A benchmark environment for \emph{multi-agent} LLM security with heterogeneous
roles, partial observability, and forced-coordination tasks across star, chain, tree,
and mesh topologies at $3$--$50$ agents---phenomena undefined in single-agent harnesses.
\item A natively parameterized adversary: adversary fraction, collusion on/off,
static/adaptive, orchestrator compromised on/off, malicious tools, Sybil identities,
and poisoned shared memory, with system-level compromise objectives.
\item A system-level metric suite generalizing ASB's $\mathrm{NRP}=\mathrm{PNA}\cdot(1-\mathrm{ASR})$
to $\NRPs=\PNAs\cdot(1-\ASRs)$, scored by a deterministic state checker so a
successful injection cannot hijack the evaluator.
\item \corb: a co-evolving red/blue evaluation protocol as a reusable methodology,
logging $(\text{attack},\text{defense})$ trajectories and their system-level scores.
\item Instrument-supporting theory: an axiomatic characterization of the aggregator,
an evaluator-integrity separation, topology-dependent compromise-propagation and
collusion bounds, a forced-coordination lower bound, and estimator concentration---all
under explicitly stated assumptions, with a headline degradation claim measured on a
transparent synthetic emulator and scoped by real-backbone replication.
\end{highlights}

\begin{keyword}
multi-agent LLM security \sep security benchmark and environment \sep co-evolving
red/blue evaluation \sep collusion \sep compromised orchestrator \sep Sybil attack
\sep prompt injection \sep tool emulation \sep deterministic security checking \sep
system-level metrics \sep compromise propagation
\end{keyword}
